\section{Introducción y contexto del problema}
\label{sec:introduccion}
 
\subsection{La reforestación y su objetivo}
La deforestación por cambio de uso de suelo, impulsada sobre todo por la expansión agrícola, la ganadería y la infraestructura, es una de las principales causas de pérdida de biodiversidad y de servicios ecosistémicos \autocite{geist2002}. Una de las respuestas a este problema es la restauración ecológica, definida como el proceso de asistir la recuperación de un ecosistema que ha sido degradado, dañado o destruido \autocite{ser2004}.
La reforestación es una de las herramientas de la restauración, la cual consiste en establecer vegetación, mediante siembra o plantación, en terrenos que antes la tenían y la perdieron \autocite{conafor2010}. Su objetivo no es solo recuperar la cobertura vegetal, sino también los servicios que el ecosistema presta, como la retención de suelo, la infiltración de agua, el hábitat para la fauna y la captura de carbono \autocite{chazdon2008}.
 
\subsection{Tipos y estrategias de reforestación}
Según el propósito, la reforestación puede ser de conservación o restauración, comercial, agroforestal o urbana \autocite{conafor2010}. Según el grado de intervención, se distingue entre la restauración pasiva, que retira el disturbio y deja que el sitio se regenere de forma natural, y la activa, la cual planta especies cuando la regeneración natural no es suficiente \autocite{holl2011}. Entre ambas existen estrategias intermedias como la nucleación, que consiste en crear pequeños núcleos (islas de árboles, perchas para aves, transposición de suelo) que aceleran la llegada de semillas y la regeneración alrededor de ellos \autocite{reis2010,corbin2012}. La elección depende del grado y tipo de disturbio y de los objetivos del proyecto \autocite{ser2004}.
 
\subsection{Cómo se realiza una reforestación}
En la mayoría de los casos, una reforestación sigue tres etapas \autocite{conafor2010}. La planeación incluye el diagnóstico del sitio (clima, suelo, pendiente, vegetación existente y disturbios), la selección de especies, la definición de la densidad y del diseño de plantación, y la obtención de la planta en vivero. La ejecución comprende la preparación del terreno, el transporte de la planta y la plantación en la temporada de lluvias. Finalmente, el mantenimiento y seguimiento abarcan la protección contra el ganado y el fuego, la reposición de plantas muertas y la medición de la supervivencia.
 
\subsection{Selección de especies, densidad y distribución}
Las especies se eligen a partir del tipo de vegetación del área que se tiene, dando prioridad a especies nativas adaptadas al clima y al suelo local \autocite{conafor2010}. Una forma práctica de hacerlo es un levantamiento de la vegetación existente en la región para plantar las especies en proporción a su frecuencia de aparición.
La densidad, es decir, el número de plantas por hectárea, y su distribución dependen tanto del tamaño que alcanza cada especie en edad adulta, como de la disponibilidad de agua y nutrientes, de la pendiente y del objetivo de la plantación \autocite{conafor2010}. Los diseños más usados son el marco real (cuadrícula) y el tresbolillo, en el que las filas se desfasan para formar triángulos equiláteros. El tresbolillo aprovecha mejor el espacio y protege mejor contra el viento y la erosión.
 
\subsection{Competencia entre plantas}
La competencia ocurre cuando plantas vecinas usan los mismos recursos limitados (agua, luz y nutrientes del suelo) y el consumo de una reduce lo disponible para la otra planta \autocite{craine2013}. Es más intensa entre plantas que están cerca, que son de tamaño parecido y que toman los recursos de la misma forma, a la misma profundidad de raíz y en la misma época.
Sin embargo, en zonas áridas las interacciones no son solo de competencia. Algunas plantas, conocidas como nodrizas, crean bajo su copa un microambiente con más sombra, humedad y nutrientes que favorece el establecimiento de otras especies \autocite{flores2003}. El mezquite (\textit{Prosopis}) es una nodriza común en el norte de México, y aprovechar este efecto es una estrategia reconocida para restaurar ambientes degradados \autocite{padilla2006}.
 
\subsection{Especies invasoras y monocultivos}
Una especie invasora es una especie exótica que se establece, se reproduce y se dispersa fuera de su área natural. Puede desplazar a las especies nativas, alterar el régimen de fuego, el ciclo de nutrientes y el uso del agua, y su control suele ser muy costoso \autocite{mack2000}. Introducirla en una reforestación, por ejemplo al usar especies de rápido crecimiento ajenas al sitio, puede dañar más el ecosistema de lo que se pretendía restaurar.
Un monocultivo es una plantación de una sola especie. Aunque es más sencillo de manejar, reduce la biodiversidad, hace a la plantación más vulnerable a plagas, enfermedades y eventos climáticos extremos y presta menos servicios ecosistémicos que una plantación mixta \autocite{liu2018}.
 
\subsection{Disponibilidad de agua}
En zonas áridas y semiáridas el agua es el recurso más limitante, por lo que condiciona tanto la selección como la distribución de las especies. Conviene elegir especies tolerantes a la sequía, como agaves, nopales y yucas, y plantarlas al inicio de la temporada de lluvias. En la distribución, plantar juntas especies con alta demanda de agua aumenta la competencia, mientras que colocar especies sensibles bajo nodrizas mejora su supervivencia \autocite{padilla2006}.
 
\subsection{Problemas de una mala planeación}
Una reforestación mal planeada o mal ejecutada puede introducir especies invasoras, generar monocultivos o acomodar las plantas de forma que compitan entre sí, lo que provoca baja supervivencia y deja sin cumplir el objetivo inicial \autocite{conafor2010,holl2011}. A esto se suma un problema logístico: si se compran plantas de más se gasta dinero de forma innecesaria y muchas mueren, y si se compran de menos no se alcanza la meta de plantación.
 
\subsection{Introducción al problema del reto}
El reto se ubica en la zona de la Línea de Transmisión Eléctrica Dominica--Charcas, en San Luis Potosí, dentro de la subprovincia Llanos y Sierras Potosino--Zacatecanos. La vegetación es principalmente matorral desértico rosetófilo y micrófilo, con un deterioro notable por sobrepastoreo de ganado. El socio formador (CONAFOR) planea reforestar 30 polígonos de entre 1.28 y 8~ha con 10 especies nativas de agaves, nopales, mezquite y yuca, en un diseño de tresbolillo elegido por los fuertes vientos de la zona y con una meta de 658 plantas por hectárea.
El problema tiene dos partes. La primera es decidir cómo distribuir las especies dentro de cada polígono para reducir la competencia entre plantas vecinas, respetando la cantidad requerida por especie. La segunda es estimar cuántas plantas de cada especie se deben comprar, considerando que en cada polígono ya existen plantas que nacieron de forma natural y cuyo número no se conoce con certeza.
 
\subsection{Justificación}
Una buena distribución de especies reduce la competencia y aprovecha la facilitación entre plantas, lo que aumenta la supervivencia y con ello el éxito de la reforestación. Esto es especialmente importante en una zona árida, donde cada planta perdida es difícil de reponer. Por otro lado, estimar bien la cantidad a comprar evita tanto el gasto innecesario como el incumplimiento de la meta. Contar con un método sistemático para ambas decisiones permite a CONAFOR planear con datos en lugar de depender de un orden de plantación fijo.
 
\subsection{Objetivo general}
El objetivo es proponer un modelo matemático y computacional que permita a CONAFOR distribuir las especies dentro de cada polígono minimizando la competencia entre ellas y cumpliendo la demanda por especie, y estimar mediante simulación la cantidad de plantas por especie que conviene adquirir considerando la incertidumbre sobre las plantas ya existentes.